\documentclass[11pt]{amsart}
\usepackage[margin=1.05in]{geometry}
\usepackage{amsmath,amssymb,amsthm,mathtools}
\usepackage{booktabs,longtable,array}
\usepackage[hidelinks]{hyperref}
\hypersetup{pdftitle={Erdos's conjecture on sum f(n)/n is false},pdfauthor={Keio}}
\DeclareTextFontCommand{\texttt}{\ttfamily\spaceskip=.5em plus .33em minus .1em\xspaceskip=\spaceskip\relax}

\newtheorem{theorem}{Theorem}[section]
\newtheorem{lemma}[theorem]{Lemma}
\newtheorem{proposition}[theorem]{Proposition}
\newtheorem{corollary}[theorem]{Corollary}
\newtheorem{conjecture}[theorem]{Conjecture}
\newtheorem*{mainA}{Theorem A}
\newtheorem*{mainB}{Theorem B}
\newtheorem*{mainC}{Theorem C}
\newtheorem*{mainD}{Theorem D}
\newtheorem*{mainE}{Theorem E}
\theoremstyle{definition}
\newtheorem{definition}[theorem]{Definition}
\newtheorem{problem}[theorem]{Problem}
\theoremstyle{remark}
\newtheorem{remark}[theorem]{Remark}

\providecommand{\C}{\mathbb C}
\providecommand{\R}{\mathbb R}
\providecommand{\Q}{\mathbb Q}
\providecommand{\Z}{\mathbb Z}
\providecommand{\N}{\mathbb N}
\providecommand{\Qb}{\overline{\mathbb Q}}
\newcommand{\Zt}{\mathbb Z_{(2)}}
\newcommand{\F}{\mathbb F}
\newcommand{\Ok}{\mathcal O}
\newcommand{\ip}[2]{\langle #1,#2\rangle}
\DeclareMathOperator{\rad}{rad}
\DeclareMathOperator{\ind}{ind}
\numberwithin{equation}{section}

\begin{document}
\title[Erd\H{o}s's conjecture on $\sum f(n)/n$ is false]{Erd\H{o}s's conjecture on $\sum f(n)/n$ is false:\\ a $2$-adic shadow of Okada's criterion and the least counterexample $q=105^2$}
\author{Keio}
\date{24 October 2026}

\begin{abstract}
Let $f$ be periodic modulo $q$ with $f(n)=\pm1$ for $q\nmid n$ and $f(n)=0$ for $q\mid n$.
Erd\H{o}s conjectured, in a letter to Livingston, that $\sum_{n\ge1}f(n)/n\neq0$ whenever the series converges.
By the theorem of Baker, Birch and Wirsing such a sum is either $0$ or transcendental, so the conjecture is a transcendence statement; it had been proved for prime $q$, for $2\varphi(q)+1>q$ (Okada), for $q\equiv3\pmod 4$ (Murty--Saradha), and for a set of $q\equiv1\pmod4$ of lower density at least $0.82$ (Chatterjee--Murty).
We show that the conjecture is false.
For $q=11025=3^2 5^2 7^2$ we give an explicit $f$, built from Legendre symbols on the divisor layers of $\Z/q\Z$, with $\sum f(n)/n=0$; the proof reduces to eleven rational identities that can be checked by hand.
Further counterexamples exist for $q=165^2$, $315^2$, $495^2$ and $735^2$, and the set of moduli carrying counterexamples is closed under multiplication.

The counterexamples come from a $2$-adic principle.
Since $\pm1\equiv1\pmod 2$, an Erd\H{o}s function with vanishing sum is an \emph{odd vector} of the Okada lattice, and we show that the existence of odd vectors is decided after a Galois descent to the maximal $2$-power quotient of $(\Z/q\Z)^\times$ (the ``$2$-shadow'').
Every odd square carries odd vectors.
For every odd non-square $q\le13105$ there are none (exact computation with explicit certificates); this proves Erd\H{o}s's conjecture, and hence the transcendence of all the sums, for these $q$, including $74$ moduli not covered by earlier criteria.
Consequently, for $q\le 13105$ the conjecture fails exactly when $q=11025$.
We conjecture that the $2$-adic obstruction holds for every odd non-square, which would confine all counterexamples to square moduli.
\end{abstract}
\maketitle
\tableofcontents

\section{Introduction}\label{sec:intro}

\subsection{The conjecture}
For an integer $q\ge2$ call a function $f\colon\Z\to\{-1,0,1\}$ an \emph{Erd\H{o}s function modulo $q$} if $f$ is $q$-periodic, $f(n)=\pm1$ for $q\nmid n$ and $f(n)=0$ for $q\mid n$.
Put
\[
L(1,f)=\sum_{n\ge1}\frac{f(n)}{n}.
\]
The series converges if and only if $\sum_{n=1}^{q}f(n)=0$; since this is a sum of $q-1$ terms $\pm1$, convergence forces $q$ to be odd.

\begin{conjecture}[Erd\H{o}s; see Livingston \cite{Liv65}]\label{conj:erdos}
If $f$ is an Erd\H{o}s function modulo $q$ and $\sum_{n=1}^{q}f(n)=0$, then $L(1,f)\neq0$.
\end{conjecture}

The problem belongs to the circle of questions raised by Chowla \cite{Cho70} on the non-vanishing of $L(1,f)$ for periodic $f$.
Baker, Birch and Wirsing \cite{BBW73} answered Chowla's question for prime moduli, which settles Conjecture~\ref{conj:erdos} for prime $q$.
For algebraic-valued periodic $f$ the number $L(1,f)$ is, by Baker's theorem on linear forms in logarithms \cite{Bak75}, either $0$ or transcendental (see \cite{BBW73,MS07}; for related transcendence results see \cite{ST03}); thus Conjecture~\ref{conj:erdos} is equivalent to the statement that \emph{every convergent Erd\H{o}s sum $\sum f(n)/n$ is transcendental}.
Okada \cite{Oka82} gave a criterion for $L(1,f)=0$ as a system of $\varphi(q)+\omega(q)$ rational linear equations and proved the conjecture when $2\varphi(q)+1>q$.
Murty and Saradha \cite{MS10} proved it for $q\equiv 3\pmod 4$.
Chatterjee and Murty \cite{CM15} proved it for a set of $q\equiv1\pmod4$ of lower density at least $0.82$, and Pathak \cite{Pat19} showed that it holds ``with probability one'' over Erd\H{o}s functions.
Murty and Pathak \cite{MP22} found an explicit basis of the Okada space and remarked that it ``does not seem to help'' for $q\equiv1\pmod4$; Bharadwaj \cite{Bha22} recast Okada's criterion through primitive characters and obtained a further family of admissible $q$.
Livingston's own proposed route was shown to fail for composite $q$ by Pathak \cite{Pat17}.
As far as we are aware the conjecture has been regarded as open for $q\equiv1\pmod4$ (it is stated as open in \cite{Pat19,MP22,Bha22}).

\subsection{Results}
\begin{mainA}[Theorem~\ref{thm:11025}]
Let $q=11025=3^2\cdot5^2\cdot7^2$.
There is an Erd\H{o}s function $f$ modulo $q$, given explicitly in Table~\ref{tab:11025} in terms of Legendre symbols, with $\sum_{n=1}^{q}f(n)=0$ and
\[
\sum_{n\ge1}\frac{f(n)}{n}=0 .
\]
In particular Conjecture~\ref{conj:erdos} is false.
\end{mainA}

The proof of Theorem~A is elementary: by a decomposition of $L(1,f)$ into Dirichlet $L$-values and logarithms of primes (Proposition~\ref{prop:decomp}), the vanishing of $L(1,f)$ follows from seven identities between rational numbers attached to the seven non-trivial real characters of conductor dividing $105$, three identities attached to $\log 3,\log5,\log7$, and the convergence identity; all eleven are listed in Tables~\ref{tab:chars} and~\ref{tab:logs}.
High-precision numerics give $|L(1,f)|<10^{-130}$ as an independent check.

\begin{mainB}[Theorems~\ref{thm:27225} and~\ref{thm:semigroup}]
There is also an Erd\H{o}s function modulo $165^2=27225$ with vanishing sum.
If $q$ and $N$ carry Erd\H{o}s functions with vanishing sums, so does $qN$.
Hence every $q=11025^a\,27225^b$ with $a+b\ge1$ carries one; so do $315^2$, $495^2$ and $735^2$ (Theorem~\ref{thm:more}).
\end{mainB}

The counterexamples were not found by blind search: the space of Erd\H{o}s functions modulo $11025$ has $2^{11024}$ elements.
They come from a $2$-adic principle which we now describe.
Because $\pm1\equiv1\pmod2$, an Erd\H{o}s function with $L(1,f)=0$ is a vector of the Okada lattice all of whose coordinates are odd; we call such vectors \emph{odd Okada vectors} (Definition~\ref{def:odd}).
Their existence is a question about the $2$-adic integral structure of the transcendental relations among the numbers $\psi(a/q)$, and it can be settled by finite linear algebra.

\begin{mainC}[$2$-shadow descent, Theorem~\ref{thm:shadow}]
Let $H\le(\Z/q\Z)^\times$ be the subgroup of elements of odd order.
If $q$ has an odd Okada vector, it has one that is invariant under multiplication by $H$.
For such vectors Okada's criterion only involves characters of $2$-power order, and it becomes a linear system on the maximal $2$-power quotients $\prod_{p\mid d}C_{2^{v_2(p-1)}}$ of the groups $(\Z/d\Z)^\times$, $d\mid q$, with $\prod_{p^e\parallel q}(1+e\,2^{v_2(p-1)})-1$ unknowns.
\end{mainC}

\begin{mainD}[Theorems~\ref{thm:squares}, \ref{thm:nonsquares} and Corollary~\ref{cor:transc}]
\begin{enumerate}
\item Every odd square has an odd Okada vector.
\item No odd non-square $q\le 13105$ has an odd Okada vector. Consequently Conjecture~\ref{conj:erdos} holds for every odd non-square $q\le13105$, and for these $q$ every convergent Erd\H{o}s sum is transcendental. Among them are $74$ moduli $q\equiv1\pmod 4$ (the smallest are $105$, $165$, $525$, $585$, $825$ and $945$) not covered by the criteria of \cite{BBW73,Oka82,MS10,CM15,Bha22}.
\end{enumerate}
\end{mainD}

Part (2) is a computer-assisted proof: for each of the $4938$ odd composite non-squares $q\le13105$ we exhibit a rational combination of the shadow equations whose coefficients are $2$-integral and whose value on the all-ones vector is odd (Section~\ref{sec:nonsquares} and Appendix~\ref{app:cert}).
Together with Theorem~A and Okada's theorem this gives a complete picture in a long initial range.

\begin{mainE}[Theorem~\ref{thm:least}]
For $2\le q\le13105$, an Erd\H{o}s function modulo $q$ with convergent and vanishing sum exists if and only if $q=11025$.
In particular $11025$ is the least modulus for which Conjecture~\ref{conj:erdos} fails.
\end{mainE}

We expect part~(2) of Theorem~D to hold for all odd non-squares (Conjecture~\ref{conj:nonsquare}); this would prove Conjecture~\ref{conj:erdos} for all non-square $q$ and confine all counterexamples to square moduli.
The shadow system also gives unconditional results for infinite families: a residuacity criterion (Theorem~\ref{thm:resid}) proves Conjecture~\ref{conj:erdos} whenever some prime $p\parallel q$ has all other prime factors of $q$ as $2^{v_2(p-1)}$-th power residues, which yields infinitely many new $q\equiv1\pmod4$ (for example $q=3\cdot7\cdot13\cdot19\cdot31$); and the same principle gives a four-line proof of the theorem of Murty and Saradha for $q\equiv3\pmod4$ (Proposition~\ref{prop:MS}).

\subsection{The idea in one paragraph}
Okada's criterion expresses $L(1,f)=0$ as a rational linear system; its solution space $\Ok_q$ is stable under multiplication of the argument by units.
An Erd\H{o}s function with $L(1,f)=0$ lies in $\Ok_q$ and is $\equiv1$ modulo $2$ at every non-zero residue.
Averaging such a vector over the odd-order part $H$ of the unit group keeps it $2$-integral and odd, because $|H|$ is odd, so nothing $2$-adic is lost by passing to $H$-invariant vectors (Theorem~C).
The remaining system lives on small $2$-groups and can be decided exactly.
When it has no odd solution (which, as far as we can test, happens precisely for non-squares) the conjecture follows for $q$.
When it does have odd solutions (squares), one may look for solutions with values $\pm1$; lifting such a shadow solution back along the fibres of the projection gives an honest Erd\H{o}s function with $L(1,f)=0$.
This is how the counterexamples of Theorems~A and~B were found (Section~\ref{sec:found}).

\subsection{Organization}
Section~\ref{sec:okada} recalls Okada's criterion in a layered form and proves the decomposition used for Theorem~A.
Section~\ref{sec:counter} contains the counterexamples and the semigroup property.
Section~\ref{sec:twoadic} introduces odd Okada vectors, re-proves the Murty--Saradha theorem, and proves the shadow descent.
Section~\ref{sec:nonsquares} treats squares and non-squares and proves Theorem~E.
Section~\ref{sec:found} explains the search, and Section~\ref{sec:open} lists open problems.
The appendices contain the second counterexample, the certificate format, and a description of the accompanying code and data.

\section{Okada's criterion in layered form}\label{sec:okada}

Throughout, $q\ge3$ is odd, $q=\prod_{p\mid q}p^{e_p}$, $U_d=(\Z/d\Z)^\times$, and $\varphi$ is Euler's function.
For a divisor $d>1$ of $q$ the \emph{layer} of level $d$ is the set of residues $n\bmod q$ with $\gcd(n,q)=q/d$, that is $n=(q/d)x$ with $x\in U_d$.
The non-zero residues modulo $q$ are the disjoint union of the layers, and $\sum_{d\mid q,\,d>1}\varphi(d)=q-1$.
For a $q$-periodic $f$ with $f(0)=0$ we write
\[
f_d(x)=f\bigl((q/d)x\bigr)\quad(x\in U_d),\qquad S_d=\sum_{x\in U_d}f_d(x),\qquad \ip{f_d}{\chi}=\sum_{x\in U_d}f_d(x)\overline{\chi(x)}
\]
for Dirichlet characters $\chi$ modulo $d$.
A character $\chi$ modulo $d$ is induced by a unique primitive character $\chi^*$ whose conductor $f_\chi$ divides $d$.
For a prime $p\mid q$ and a layer $d$ put
\begin{equation}\label{eq:eps}
\varepsilon_p(d)=v_p(q/d)-\frac{[\,p\mid d\,]}{p-1}.
\end{equation}

\begin{proposition}[Decomposition]\label{prop:decomp}
Let $f$ be rational-valued, $q$-periodic, with $f(0)=0$ and $\sum_{n=1}^{q}f(n)=\sum_{d}S_d=0$.
Then
\begin{equation}\label{eq:decomp}
L(1,f)=\sum_{\chi^*}A_{\chi^*}(f)\,L(1,\chi^*)\;-\;\frac1q\sum_{p\mid q}\Bigl(\sum_{d\mid q,\,d>1}S_d\,\varepsilon_p(d)\Bigr)\log p,
\end{equation}
where $\chi^*$ runs over the non-principal primitive characters whose conductor divides $q$, and
\begin{equation}\label{eq:A}
A_{\chi^*}(f)=\frac1q\sum_{\substack{d\mid q\\ f_{\chi}\mid d}}\frac{d}{\varphi(d)}\prod_{\substack{p\mid d\\ p\nmid f_\chi}}\Bigl(1-\frac{\chi^*(p)}{p}\Bigr)\ip{f_d}{\chi^*}.
\end{equation}
\end{proposition}

\begin{proof}
For $s>1$ group the terms of $L(s,f)=\sum f(n)n^{-s}$ according to $\gcd(n,q)=q/d$; writing $n=(q/d)x$ with $\gcd(x,d)=1$ and expanding $f_d$ in characters modulo $d$,
\[
L(s,f)=\sum_{d}\Bigl(\frac dq\Bigr)^{s}\sum_{\gcd(x,d)=1}\frac{f_d(x)}{x^{s}}
=\sum_d\Bigl(\frac dq\Bigr)^{s}\frac1{\varphi(d)}\sum_{\chi\bmod d}\ip{f_d}{\chi}L(s,\chi).
\]
For non-principal $\chi$, $L(s,\chi)=L(s,\chi^*)\prod_{p\mid d,\,p\nmid f_\chi}(1-\chi^*(p)p^{-s})$ is continuous at $s=1$, and collecting the terms with a given $\chi^*$ gives the first sum in \eqref{eq:decomp}.
For the principal character, $L(s,\chi_0)=\zeta(s)\prod_{p\mid d}(1-p^{-s})$, and as $s\to1^+$
\[
\Bigl(\frac dq\Bigr)^{s}\frac{S_d}{\varphi(d)}\zeta(s)\prod_{p\mid d}(1-p^{-s})
=\frac{S_d}{q}\Bigl(\frac1{s-1}+\gamma-\log\frac qd+\sum_{p\mid d}\frac{\log p}{p-1}+O(s-1)\Bigr).
\]
Summing over $d$, the pole and the $\gamma$-terms cancel because $\sum_dS_d=0$, and what remains is $-\frac1q\sum_dS_d\sum_{p}\varepsilon_p(d)\log p$.
Since the Dirichlet series converges at $s=1$, $L(1,f)=\lim_{s\to1^+}L(s,f)$ by Abel's theorem.
\end{proof}

The coefficient of $\log p$ can also be written $-\frac1q\sum_{n=1}^{q-1}f(n)\varepsilon_p(n)$ with $\varepsilon_p(n)=v_p(\gcd(n,q))-[v_p(\gcd(n,q))<e_p]/(p-1)$.
Proposition~\ref{prop:decomp} immediately gives the sufficiency half of the following criterion; necessity rests on Baker's theorem.

\begin{theorem}[Okada \cite{Oka82}; layered form, cf.\ Bharadwaj {\cite[Thm.~4]{Bha22}}]\label{thm:okada}
Let $f$ be as in Proposition~\ref{prop:decomp}.
Then $L(1,f)=0$ if and only if
\begin{enumerate}
\item[(A)] $A_{\chi^*}(f)=0$ for every non-principal primitive $\chi^*$ of conductor dividing $q$, and
\item[(B)] $\sum_{d}S_d\,\varepsilon_p(d)=0$ for every prime $p\mid q$.
\end{enumerate}
\end{theorem}

\begin{proof}[Sketch of necessity]
For odd $\chi^*$ one has $L(1,\chi^*)\in\pi i\,\Qb$, and for even $\chi^*\neq1$ one has $L(1,\chi^*)=-\frac{\tau(\chi^*)}{f_\chi}\sum_{a}\overline{\chi^*}(a)\log|1-\zeta_{f_\chi}^{a}|$.
The number $\pi i=\log(-1)$ is $\Q$-linearly independent of logarithms of positive real algebraic numbers, so by Baker's theorem the $\pi$-part and the real-logarithmic part of \eqref{eq:decomp} vanish separately.
By Gauss's formula for $\psi(a/q)$ the $\pi$-part is $\frac{\pi}{2q}\sum_af(a)\cot(\pi a/q)=-\frac{\pi i}{2q}X$ with $X=\sum_af(a)\frac{1+\zeta_q^a}{1-\zeta_q^a}\in\Q(\zeta_q)$ (here $f$ is rational), and $X=0$ if and only if all its Galois conjugates $\sum_af(a)\frac{1+\zeta_q^{ca}}{1-\zeta_q^{ca}}$, $c\in U_q$, vanish; since replacing $f$ by $f(c^{-1}\cdot)$ multiplies $A_{\chi^*}$ by $\overline{\chi^*}(c)$, orthogonality and $L(1,\chi^*)\neq0$ give $A_{\chi^*}=0$ for odd $\chi^*$.
For the logarithmic part one decomposes the $\Qb$-span of the classes of $\log|1-\zeta^a|$ and $\log p$ into eigenspaces of $\mathrm{Gal}(\Q(\zeta_q)/\Q)$: the even $L(1,\chi^*)$ are non-zero vectors in pairwise distinct eigenspaces, and the classes of $\log p$, $p\mid q$, are linearly independent in the invariant eigenspace; Baker's theorem turns this into linear independence of the corresponding numbers over $\Qb$.
This is the argument of \cite{BBW73,Oka82}; in the present layered form it is \cite[Thm.~4]{Bha22}.
\end{proof}

Theorem~\ref{thm:okada}, together with $\sum_dS_d=0$, is a system of $\varphi(q)+\omega(q)$ independent rational linear equations in the $q-1$ values $f(1),\dots,f(q-1)$ (the characters contribute $\varphi(q)-1$ real equations, (B) contributes $\omega(q)$).
It is equivalent to Okada's original system \cite{Oka82}: $\sum_{m\in M(q)}f(am)/m=0$ for all $a\in U_q$, where $M(q)$ is the multiplicative monoid generated by the primes dividing $q$, together with one equation for each $p\mid q$.
We denote by
\[
\Ok_q=\Bigl\{f\colon(\Z/q\Z)\smallsetminus\{0\}\to\Q\ :\ \textstyle\sum_nf(n)=0\text{ and (A), (B) hold}\Bigr\}
\]
the solution space (with $f(0)=0$); $\dim\Ok_q=q-1-\varphi(q)-\omega(q)$.
It is stable under $f\mapsto f(u\,\cdot)$ for $u\in U_q$, because this multiplies each $A_{\chi^*}(f)$ by a root of unity and leaves all $S_d$ unchanged.

\section{Counterexamples}\label{sec:counter}

\subsection{\texorpdfstring{The modulus $105^2$}{The modulus 105^2}}
Let $q=11025=3^2\cdot5^2\cdot7^2$.
For $n\not\equiv0\pmod q$ write $n=(q/d)x$ with $d\mid q$, $d>1$, $x\in U_d$, and for $p\mid d$ put $\lambda_p=\bigl(\frac xp\bigr)$ (Legendre symbol).
Define $f(n)=F_d(x)$ with $F_d$ as in Table~\ref{tab:11025} and $f(n)=0$ for $q\mid n$.
Each $F_d$ is a $\pm1$-valued function of the Legendre symbols $\lambda_p$, $p\mid d$, so $f$ is an Erd\H{o}s function modulo $q$.
For $m\mid105$ let $\chi_m$ be the real primitive character of conductor $m$, $\chi_m(x)=\prod_{p\mid m}\bigl(\frac xp\bigr)$, and write $\lambda_m=\chi_m(x)$.
The last column of Table~\ref{tab:11025} records the expansion $F_d=\sum_{m\mid\rad(d)}c_{d,m}\lambda_m$.

\begin{longtable}{rrrll}
\caption{The counterexample modulo $q=11025$: for $n$ with $\gcd(n,q)=q/d$ write $n=(q/d)x$; then $f(n)=F_d(x)$, where $\lambda_p=\bigl(\tfrac{x}{p}\bigr)$. The last column gives the expansion $F_d=\sum_m c_{d,m}\lambda_m$ ($\lambda_m=\prod_{p\mid m}\lambda_p$).}\label{tab:11025}\\
\toprule $d$ & $q/d$ & $\varphi(d)$ & $F_d(x)$ & nonzero $c_{d,m}$\\ \midrule\endfirsthead
\toprule $d$ & $q/d$ & $\varphi(d)$ & $F_d(x)$ & nonzero $c_{d,m}$\\ \midrule\endhead
3 & 3675 & 2 & $-1$ & $c_{1}=-1$\\
5 & 2205 & 4 & $-\lambda_5$ & $c_{5}=-1$\\
7 & 1575 & 6 & $-\lambda_7$ & $c_{7}=-1$\\
9 & 1225 & 6 & $-1$ & $c_{1}=-1$\\
15 & 735 & 8 & $-\lambda_3\lambda_5$ & $c_{15}=-1$\\
21 & 525 & 12 & $1$ if $\lambda_3=1$; $-\lambda_7$ if $\lambda_3=-1$ & $c_{1}=\tfrac{1}{2}$, $c_{7}=-\tfrac{1}{2}$, $c_{3}=\tfrac{1}{2}$, $c_{21}=\tfrac{1}{2}$\\
25 & 441 & 20 & $\lambda_5$ & $c_{5}=1$\\
35 & 315 & 24 & $\lambda_5\lambda_7$ & $c_{35}=1$\\
45 & 245 & 24 & $-1$ if $\lambda_3=1$; $\lambda_5$ if $\lambda_3=-1$ & $c_{1}=-\tfrac{1}{2}$, $c_{5}=\tfrac{1}{2}$, $c_{3}=-\tfrac{1}{2}$, $c_{15}=-\tfrac{1}{2}$\\
49 & 225 & 42 & $-\lambda_7$ & $c_{7}=-1$\\
63 & 175 & 36 & $1$ if $\lambda_3=1$; $-\lambda_7$ if $\lambda_3=-1$ & $c_{1}=\tfrac{1}{2}$, $c_{7}=-\tfrac{1}{2}$, $c_{3}=\tfrac{1}{2}$, $c_{21}=\tfrac{1}{2}$\\
75 & 147 & 40 & $1$ if $\lambda_3=1$; $-\lambda_5$ if $\lambda_3=-1$ & $c_{1}=\tfrac{1}{2}$, $c_{5}=-\tfrac{1}{2}$, $c_{3}=\tfrac{1}{2}$, $c_{15}=\tfrac{1}{2}$\\
105 & 105 & 48 & $-\lambda_7$ if $\lambda_3=1$; $-1$ if $\lambda_3=-1$ & $c_{1}=-\tfrac{1}{2}$, $c_{7}=-\tfrac{1}{2}$, $c_{3}=\tfrac{1}{2}$, $c_{21}=-\tfrac{1}{2}$\\
147 & 75 & 84 & $\lambda_7$ & $c_{7}=1$\\
175 & 63 & 120 & $-\lambda_5\lambda_7$ & $c_{35}=-1$\\
225 & 49 & 120 & $\lambda_3$ & $c_{3}=1$\\
245 & 45 & 168 & $\lambda_5\lambda_7$ & $c_{35}=1$\\
315 & 35 & 144 & $\lambda_5\lambda_7$ & $c_{35}=1$\\
441 & 25 & 252 & $-\lambda_3$ & $c_{3}=-1$\\
525 & 21 & 240 & $\lambda_5\lambda_7$ & $c_{35}=1$\\
735 & 15 & 336 & $\lambda_5\lambda_7$ & $c_{35}=1$\\
1225 & 9 & 840 & $\lambda_5\lambda_7$ & $c_{35}=1$\\
1575 & 7 & 720 & $\lambda_5\lambda_7$ & $c_{35}=1$\\
2205 & 5 & 1008 & $\lambda_5\lambda_7$ & $c_{35}=1$\\
3675 & 3 & 1680 & $\lambda_5\lambda_7$ & $c_{35}=1$\\
11025 & 1 & 5040 & $-\lambda_5\lambda_7$ & $c_{35}=-1$\\
\bottomrule\end{longtable}


\begin{theorem}\label{thm:11025}
The function $f$ of Table~\ref{tab:11025} satisfies $\sum_{n=1}^{q}f(n)=0$ and
\[
\sum_{n=1}^{\infty}\frac{f(n)}{n}=0.
\]
\end{theorem}

\begin{proof}
We apply Proposition~\ref{prop:decomp}.
Since each $F_d$ depends only on Legendre symbols, $F_d(xy^2)=F_d(x)$ for all $x,y\in U_d$, hence $\ip{f_d}{\chi}=\overline{\chi(y)}^{\,2}\ip{f_d}{\chi}$ for all $y$; so $\ip{f_d}{\chi}=0$ unless $\chi^2=1$.
A real primitive character of odd conductor dividing $q$ has squarefree conductor, so it is one of $\chi_m$, $m\mid105$, $m>1$.
By orthogonality of the distinct characters $\lambda_m$ on $U_d$, $\ip{f_d}{\chi_m}=\varphi(d)\,c_{d,m}$ and $S_d=\varphi(d)\,c_{d,1}$.
Hence \eqref{eq:A} gives
\[
q\,A_{\chi_m}(f)=\sum_{m\mid d\mid q}w_{d,m}\,c_{d,m},\qquad w_{d,m}=d\prod_{\substack{p\mid d\\p\nmid m}}\Bigl(1-\frac{\chi_m(p)}{p}\Bigr),
\]
and all other coefficients $A_{\chi^*}(f)$ vanish.
The values $\chi_m(p)$ follow from $\bigl(\frac35\bigr)=\bigl(\frac37\bigr)=\bigl(\frac53\bigr)=\bigl(\frac57\bigr)=\bigl(\frac75\bigr)=-1$ and $\bigl(\frac73\bigr)=1$.
Table~\ref{tab:chars} lists every non-zero term $w_{d,m}c_{d,m}$; each row sums to $0$ (for $m=105$ there is no term).
Table~\ref{tab:logs} lists the non-zero layer sums $S_d$; they sum to $0$, which is the convergence condition, and $\sum_dS_d\varepsilon_p(d)=0$ for $p=3,5,7$.
By \eqref{eq:decomp}, $L(1,f)=0$.
\end{proof}

\begin{table}[ht]\centering\small
\caption{The seven identities $\sum_{m\mid d} w_{d,m}c_{d,m}=0$, with $w_{d,m}=d\prod_{p\mid d,\,p\nmid m}\bigl(1-\chi_m(p)/p\bigr)$; entries are $w_{d,m}c_{d,m}$ for the layers with $c_{d,m}\neq0$.}\label{tab:chars}
\begin{tabular}{@{}l>{\raggedright}p{0.27\textwidth}>{\raggedright\arraybackslash}p{0.62\textwidth}@{}}\toprule $m$ & $\chi_m(p)$, $p\nmid m$ & terms $w_{d,m}c_{d,m}$ (layer $d$ in brackets) \\ \midrule
3 & $\chi_{3}(5)=-1$, $\chi_{3}(7)=+1$ & $9\,[21]$ $-27\,[45]$ $27\,[63]$ $45\,[75]$ $54\,[105]$ $270\,[225]$ $-378\,[441]$ \\
5 & $\chi_{5}(3)=-1$, $\chi_{5}(7)=-1$ & $-5\,[5]$ $25\,[25]$ $30\,[45]$ $-50\,[75]$ \\
7 & $\chi_{7}(3)=-1$, $\chi_{7}(5)=-1$ & $-7\,[7]$ $-14\,[21]$ $-49\,[49]$ $-42\,[63]$ $-84\,[105]$ $196\,[147]$ \\
15 & $\chi_{15}(7)=-1$ & $-15\,[15]$ $-\tfrac{45}{2}\,[45]$ $\tfrac{75}{2}\,[75]$ \\
21 & $\chi_{21}(5)=+1$ & $\tfrac{21}{2}\,[21]$ $\tfrac{63}{2}\,[63]$ $-42\,[105]$ \\
35 & $\chi_{35}(3)=+1$ & $35\,[35]$ $-175\,[175]$ $245\,[245]$ $210\,[315]$ $350\,[525]$ $490\,[735]$ $1225\,[1225]$ $1050\,[1575]$ $1470\,[2205]$ $2450\,[3675]$ $-7350\,[11025]$ \\
105 & --- &  \\
\bottomrule\end{tabular}\end{table}

\begin{table}[ht]\centering\small
\caption{Layer sums $S_d=\varphi(d)c_{d,1}$ (layers with $S_d\ne0$) and the weights $\varepsilon_p(d)=v_p(q/d)-[p\mid d]/(p-1)$.}\label{tab:logs}
\begin{tabular}{rrrrr}\toprule $d$ & $S_d$ & $\varepsilon_3(d)$ & $\varepsilon_5(d)$ & $\varepsilon_7(d)$ \\ \midrule
3 & -2 & $\tfrac{1}{2}$ & $2$ & $2$ \\
9 & -6 & $-\tfrac{1}{2}$ & $2$ & $2$ \\
21 & 6 & $\tfrac{1}{2}$ & $2$ & $\tfrac{5}{6}$ \\
45 & -12 & $-\tfrac{1}{2}$ & $\tfrac{3}{4}$ & $2$ \\
63 & 18 & $-\tfrac{1}{2}$ & $2$ & $\tfrac{5}{6}$ \\
75 & 20 & $\tfrac{1}{2}$ & $-\tfrac{1}{4}$ & $2$ \\
105 & -24 & $\tfrac{1}{2}$ & $\tfrac{3}{4}$ & $\tfrac{5}{6}$ \\
\midrule $\sum$ & 0 & $\sum S_d\varepsilon_3=0$ & $\sum S_d\varepsilon_5=0$ & $\sum S_d\varepsilon_7=0$ \\
\bottomrule\end{tabular}\end{table}


\begin{remark}\label{rem:11025checks}
The theorem was checked in three further ways:
(i) all $5040$ of Okada's unit equations $\sum_{m\in M(q)}f(am)/m=0$ and his three prime equations hold exactly in rational arithmetic, with the infinite sums evaluated as finite combinations of geometric series;
(ii) evaluating $L(1,f)=-\frac1q\sum_{a=1}^{q-1}f(a)\psi(a/q)$ in PARI/GP \cite{PARI} with $134$ significant digits gives $L(1,f)\approx-1.2\cdot10^{-135}$, i.e.\ zero to working precision, while exchanging the two values $f(1)=-1$ and $f(2)=1$ changes it to $1.00000003\ldots$;
(iii) the partial sum $\sum_{n\le 5q}f(n)/n\approx2.3\cdot10^{-5}$ is of the size expected for a convergent series with value $0$.
\end{remark}

\subsection{\texorpdfstring{The modulus $165^2$}{The modulus 165^2}}
\begin{theorem}\label{thm:27225}
Let $q=27225=3^2\cdot 5^2\cdot 11^2$ and let $f$ be the Erd\H{o}s function modulo $q$ given in Table~\ref{tab:27225} (Appendix~\ref{app:27225}).
Then $\sum_{n=1}^{q}f(n)=0$ and $L(1,f)=0$.
\end{theorem}

\begin{proof}[Proof (computer-verified)]
The function $f$ is invariant under the subgroup $H=\{u^{16}:u\in U_q\}$ of odd-order units ($|H|=825$), so $\ip{f_d}{\chi}=0$ unless $\chi$ has $2$-power order.
The primitive characters of $2$-power order and conductor dividing $q$ are $\chi=\bigl(\frac{\cdot}{3}\bigr)^{a}\,\omega_5^{\,b}\,\bigl(\frac{\cdot}{11}\bigr)^{c}$ with $a,c\in\{0,1\}$, $b\in\{0,1,2,3\}$, where $\omega_5$ is the quartic character modulo $5$ with $\omega_5(2)=i$.
For each of the $15$ non-trivial ones the coefficient \eqref{eq:A} was computed exactly in $\Q(i)$ and found to be $0$; the three sums $\sum_dS_d\varepsilon_p(d)$, $p\in\{3,5,11\}$, and $\sum_dS_d$ are $0$.
As independent checks, all $13200$ unit equations and the three prime equations of Okada's original system hold exactly, and PARI/GP gives $L(1,f)\approx-1.1\cdot10^{-135}$ with $134$ significant digits.
\end{proof}

Unlike the first example, no counterexample modulo $27225$ that is invariant under $H$ depends only on Legendre symbols (the corresponding finite search space is empty, Section~\ref{sec:found}); the quartic character modulo $5$ is needed.

\subsection{\texorpdfstring{Further examples: $315^2$, $495^2$ and $735^2$}{Further examples}}
Restricting the search of Section~\ref{sec:found} to $H$-invariant functions that depend only on Legendre symbols also produced counterexamples modulo
\[
315^2=3^4\cdot5^2\cdot7^2=99225,\qquad 495^2=3^4\cdot5^2\cdot11^2=245025,\qquad 735^2=3^2\cdot5^2\cdot7^4=540225 .
\]
\begin{theorem}\label{thm:more}
For $q\in\{315^2,495^2,735^2\}$ there is an Erd\H{o}s function modulo $q$ with vanishing sum which depends, on every layer, only on the Legendre symbols $\bigl(\frac xp\bigr)$, $p\mid q$.
\end{theorem}
\begin{proof}
For a function of this type the proof of Theorem~\ref{thm:11025} applies verbatim: only the real characters $\chi_m$, $m\mid\rad(q)$, can occur, and $L(1,f)=0$ follows from the seven identities $\sum_{m\mid d}w_{d,m}c_{d,m}=0$, the three logarithmic identities and the convergence identity, now summed over the $44$ layers of $q$.
For the three functions in the supplementary material all eleven identities hold exactly (script \texttt{verify\_legendre\_type.py}).
\end{proof}
For $99225$ we also checked all $45360$ unit equations and the three prime equations of Okada's original system exactly, and PARI/GP gives $|L(1,f)|\approx1.2\cdot10^{-135}$ with $134$ significant digits; for $245025$ and $540225$, PARI/GP gives $|L(1,f)|\approx7.8\cdot10^{-78}$ and $1.5\cdot10^{-77}$ with $77$ significant digits.
These examples do not follow from Theorem~\ref{thm:semigroup}, since $9$ and $49$ carry no counterexamples (Okada's theorem); note that $495^2=9\cdot165^2$ has a Legendre-type counterexample although $165^2$ has no $H$-invariant one of Legendre type.

\subsection{Products of counterexample moduli}
\begin{lemma}[Composition]\label{lem:compose}
Let $f$ be $q$-periodic and $g$ be $N$-periodic, rational-valued, with $f(0)=g(0)=0$ and $\sum_{n=1}^qf(n)=\sum_{n=1}^Ng(n)=0$.
Define the $qN$-periodic function
\[
F(n)=\begin{cases}f(n),& q\nmid n,\\ g(n/q),& q\mid n.\end{cases}
\]
Then $F(0)=0$, $\sum_{n=1}^{qN}F(n)=0$, and $L(1,F)=L(1,f)+\frac1qL(1,g)$.
If $f$ and $g$ are Erd\H{o}s functions modulo $q$ and $N$, then $F$ is an Erd\H{o}s function modulo $qN$.
\end{lemma}

\begin{proof}
Periodicity and $F(0)=g(0)=0$ are clear.
Over a period, $\sum_{q\nmid n}F(n)=N\sum_{n=1}^qf(n)=0$ and $\sum_{q\mid n}F(n)=\sum_{m=1}^Ng(m)=0$.
The series $\sum_{q\nmid n}f(n)/n=L(1,f)$ (recall $f(n)=0$ for $q\mid n$) and $\sum_{m}g(m)/(qm)=L(1,g)/q$ both converge, and their sum is $L(1,F)$.
If $f(n)=\pm1$ for $q\nmid n$ and $g(m)=\pm1$ for $N\nmid m$, then $F(n)=\pm1$ unless $q\mid n$ and $N\mid n/q$, i.e.\ unless $qN\mid n$.
\end{proof}

\begin{theorem}\label{thm:semigroup}
Let $\mathcal C$ be the set of moduli $q$ for which an Erd\H{o}s function with $\sum_{n=1}^qf(n)=0$ and $L(1,f)=0$ exists.
Then $\mathcal C$ is closed under multiplication.
In particular
\[
105^{2a}\,165^{2b}\,315^{2c}\,495^{2e}\,735^{2g}\in\mathcal C\qquad(a,b,c,e,g\ge0\text{ not all }0).
\]
\end{theorem}

\begin{proof}
Apply Lemma~\ref{lem:compose} with $L(1,f)=L(1,g)=0$, and use Theorems~\ref{thm:11025}, \ref{thm:27225} and~\ref{thm:more}.
\end{proof}

\section{\texorpdfstring{The $2$-adic principle}{The 2-adic principle}}\label{sec:twoadic}

Let $\Zt$ be the ring of rationals with odd denominator.
Every element of $\Zt$ has a residue modulo $2$, and we call it odd or even accordingly.

\begin{definition}\label{def:odd}
An \emph{odd Okada vector} for $q$ is an $f\in\Ok_q$ with $f(n)\in\Zt$ and $f(n)$ odd for every $n\not\equiv0\pmod q$.
We say that $q$ is \emph{$2$-adically obstructed} if no odd Okada vector exists.
\end{definition}

\begin{lemma}\label{lem:erdos-odd}
If $f$ is an Erd\H{o}s function modulo $q$ with $\sum_{n=1}^qf(n)=0$ and $L(1,f)=0$, then $f$ is an odd Okada vector.
In particular, Conjecture~\ref{conj:erdos} holds for every $2$-adically obstructed $q$, and for such $q$ every convergent Erd\H{o}s sum is transcendental.
\end{lemma}

\begin{proof}
By Theorem~\ref{thm:okada}, $f\in\Ok_q$, and $\pm1$ is odd.
The last statement follows from the fact that $L(1,f)$ is $0$ or transcendental \cite{BBW73,MS07}.
\end{proof}

The search for odd Okada vectors is a linear problem over the discrete valuation ring $\Zt$, and it has a dual description.

\begin{lemma}[Certificates]\label{lem:farkas}
Let $R_1,\dots,R_k\in\Q^N$ span the space of linear forms vanishing on a subspace $V\subseteq\Q^N$, and let $\mathbf 1=(1,\dots,1)$.
There is $v\in V\cap\Zt^N$ with all coordinates odd if and only if there is no $y\in\Q^k$ such that $\Lambda=\sum_iy_iR_i$ has all coordinates in $\Zt$ and $\Lambda(\mathbf 1)$ is odd.
\end{lemma}

\begin{proof}
Let $K=V\cap\Zt^N$ and $R=V^\perp\cap\Zt^N$; both are saturated $\Zt$-lattices, of ranks $\dim V$ and $N-\dim V$.
Their reductions $\overline K,\overline R\subseteq\F_2^N$ have dimensions $\dim V$ and $N-\dim V$ (a saturated lattice has a basis that reduces to an $\F_2$-basis) and are orthogonal, hence $\overline R=\overline K^{\perp}$.
An odd vector in $K$ exists iff $\mathbf 1\in\overline K=\overline R^{\perp}$, i.e.\ iff $\Lambda(\mathbf 1)$ is even for every $\Lambda\in R$.
\end{proof}

\subsection{\texorpdfstring{A four-line proof for $q\equiv3\pmod4$}{A four-line proof for q = 3 mod 4}}
The simplest instance of the principle recovers the theorem of Murty and Saradha \cite{MS10}.
Put $\zeta=e^{2\pi i/q}$ and $\rho_a=\zeta^a/(1-\zeta^a)$.
Since $q$ is odd, $1-\zeta^a$ is a unit at every prime above $2$ for $a\not\equiv0$, so $\rho_a$ is $2$-integral.

\begin{proposition}\label{prop:MS}
If $q\equiv3\pmod 4$, then $q$ is $2$-adically obstructed.
In particular Conjecture~\ref{conj:erdos} holds for $q\equiv 3\pmod 4$ \cite{MS10}.
\end{proposition}

\begin{proof}
Let $f\in\Ok_q$.
The vanishing of the $\pi$-part of \eqref{eq:decomp} (see the proof of Theorem~\ref{thm:okada}) reads
\[
\sum_{a=1}^{q-1}f(a)\cot(\pi a/q)=0;
\]
since
\[
i\cot(\pi a/q)=\frac{1+\zeta^a}{1-\zeta^a}=1+2\rho_a
\]
and $\sum_af(a)=0$, this is $\sum_af(a)\rho_a=0$.
Pairing $a$ with $-a$ and using $\rho_a+\rho_{-a}=-1$ gives $\sum_{a=1}^{q-1}\rho_a=-(q-1)/2$.
If $f$ were odd and $2$-integral, then $f(a)=1+2g(a)$ with $g(a)\in\Zt$ and
\[
0=\sum_af(a)\rho_a=-\frac{q-1}2+2\sum_ag(a)\rho_a,
\]
which is impossible when $(q-1)/2$ is odd.
\end{proof}

For $q\equiv1\pmod4$ the same identity only says $\sum_ag(a)\rho_a=(q-1)/4$, which is not contradictory by itself; one needs the interaction of the odd and even parts of Okada's system, which is what the following descent makes computable.

\subsection{\texorpdfstring{Descent to the $2$-shadow}{Descent to the 2-shadow}}
For $d\mid q$ let $H_d\le U_d$ be the subgroup of elements of odd order and $Q_d=U_d/H_d$, the maximal $2$-power quotient.
For an odd prime $p$ let $t_p=v_2(p-1)$.
Since $U_{p^e}$ is cyclic of order $p^{e-1}(p-1)$, the $2$-part of $U_{p^e}$ maps isomorphically onto that of $U_p$, and
\[
Q_d\;\cong\;\prod_{p\mid d}C_{2^{t_p}},\qquad x\mapsto\bigl(\ind_{g_p}(x\bmod p)\bmod 2^{t_p}\bigr)_{p\mid d},
\]
where $g_p$ is a primitive root modulo $p$.
In particular $Q_d$ depends only on $\rad(d)$, and $|H_d|=\varphi(d)/|Q_d|$ is odd.
The characters of $Q_d$ are exactly the characters of $2$-power order modulo $d$; their conductors are squarefree, and the primitive ones of conductor $m\mid\rad(q)$ are the characters of $Q_m$ that are non-trivial on every factor $C_{2^{t_p}}$, $p\mid m$.

\begin{theorem}[$2$-shadow descent]\label{thm:shadow}
Let $q$ be odd and $H=H_q$.
\begin{enumerate}
\item If $q$ has an odd Okada vector, then it has one that is invariant under $n\mapsto hn$ for all $h\in H$.
\item An $H$-invariant $f$ is determined by the functions $\Phi_d\colon Q_d\to\Q$, $\Phi_d(y)=\sum_{x\mapsto y}f_d(x)$ (fibre sums over $U_d\to Q_d$), and $f$ is odd and $2$-integral iff every $\Phi_d$ is.
\item For $H$-invariant $f$, $f\in\Ok_q$ if and only if the following hold.
\begin{itemize}
\item[(A$'$)] For every $m\mid\rad(q)$, $m>1$, and every primitive $\chi\in\widehat{Q_m}$:
\[
\sum_{m\mid d\mid q}\frac d{\varphi(d)}\prod_{p\mid d,\,p\nmid m}\Bigl(1-\frac{\chi(p)}p\Bigr)\sum_{y\in Q_d}\Phi_d(y)\overline{\chi(y)}=0 .
\]
\item[(B$'$)] For every $p\mid q$: $\sum_{d}\varepsilon_p(d)\sum_{y}\Phi_d(y)=0$.
\item[(C$'$)] $\sum_d\sum_y\Phi_d(y)=0$.
\end{itemize}
\end{enumerate}
Consequently, $q$ is $2$-adically obstructed if and only if the linear system \textup{(A$'$)--(C$'$)} has no solution with all $\Phi_d(y)$ odd and $2$-integral.
The number of unknowns is $\sum_{d\mid q,\,d>1}|Q_d|=\prod_{p^{e}\parallel q}\bigl(1+e\,2^{t_p}\bigr)-1$.
\end{theorem}

\begin{proof}
(1) Let $f$ be an odd Okada vector and $\bar f(n)=|H|^{-1}\sum_{h\in H}f(hn)$.
Then $\bar f\in\Ok_q$ because $\Ok_q$ is stable under multiplication by units (Section~\ref{sec:okada}); $\bar f$ is $2$-integral because $|H|$ is odd; and $\bar f(n)$ is odd, since $\sum_hf(hn)$ is a sum of an odd number of odd elements of $\Zt$.
Clearly $\bar f$ is $H$-invariant.

(2) The image of $H$ in $U_d$ is $H_d$ (a surjection of finite abelian groups maps the odd-order part onto the odd-order part), so an $H$-invariant $f$ has $f_d$ constant on the fibres of $U_d\to Q_d$, which have odd cardinality $|H_d|$; hence $f_d=\Phi_d/|H_d|$ on each fibre and the parity statement follows.

(3) For $H$-invariant $f$ and a character $\chi$ modulo $d$, $\ip{f_d}{\chi}=0$ unless $\chi$ is trivial on $H_d$, i.e.\ unless $\chi$ has $2$-power order, in which case $\ip{f_d}{\chi}=\sum_y\Phi_d(y)\overline{\chi(y)}$.
So in Theorem~\ref{thm:okada} the conditions (A) for characters not of $2$-power order hold automatically, and the remaining ones are (A$'$); (B$'$) and (C$'$) are (B) and the convergence condition, as $S_d=\sum_y\Phi_d(y)$.
The count of unknowns follows from $|Q_d|=\prod_{p\mid d}2^{t_p}$ and the fact that there are $e_p$ choices of $v_p(d)\ge1$ for each $p\mid d$.
\end{proof}

\begin{remark}\label{rem:integer-form}
For computations we use an equivalent integer form of (A$'$).
For $m\mid\rad(q)$, $m>1$, and $y\in Q_m$ put
\[
H_m(\Phi)(y)=\sum_{m\mid d\mid q}\Bigl(\prod_{p\mid q,\,p\nmid d}(p-1)\Bigr)\Bigl[\prod_{p\mid d,\,p\nmid m}\bigl(p-T_p\bigr)\,\pi_{d\to m}\Phi_d\Bigr](y),
\]
where $\pi_{d\to m}$ sums over the fibres of $Q_d\to Q_m$ and $T_p$ is translation by the image of $p$ in $Q_m$.
Multiplying (A$'$) by the non-zero constant $\frac{\varphi(m)}{m}\prod_{p\mid q,\,p\nmid m}(p-1)$ shows that (A$'$) for all primitive $\chi\in\widehat{Q_m}$ is equivalent to: $H_m(\Phi)$ has no primitive component, i.e.
\[
\sum_{\epsilon\in\{0,1\}^{\omega(m)}}(-1)^{|\epsilon|}\,H_m(\Phi)\bigl(y_\epsilon(k)\bigr)=0\qquad\text{for all }k=(k_p)\in\prod_{p\mid m}\bigl(C_{2^{t_p}}\smallsetminus\{0\}\bigr),
\]
where $y_\epsilon(k)_p=k_p$ if $\epsilon_p=0$ and $y_\epsilon(k)_p=0$ if $\epsilon_p=1$.
Together with (B$'$) multiplied by $p-1$ and (C$'$), these are integer linear forms in the unknowns $\Phi_d(y)$.
\end{remark}

\begin{remark}
The decisive feature of Theorem~\ref{thm:shadow} is that \emph{nothing $2$-adic is lost} in passing to the shadow: the odd-order part of the Galois group $\mathrm{Gal}(\Q(\zeta_q)/\Q)\cong U_q$ can be averaged away exactly because its order is odd, and what remains is governed by the $2$-power cyclotomic subfields.
The size of the problem drops from $q-1$ unknowns to $\prod(1+e_p2^{t_p})-1$; for $q=11025$ from $11024$ to $224$.
Relaxations that keep only the quadratic characters, or only characters supported on few primes, or only the odd (resp.\ even) characters, do \emph{not} suffice in general (Section~\ref{sec:found}); the full $2$-power quotient is needed.
\end{remark}

\section{Squares and non-squares}\label{sec:nonsquares}

\subsection{Squares are not obstructed}
\begin{proposition}\label{prop:p2}
Let $p$ be an odd prime, $q=p^2$, and $\chi_p=\bigl(\frac{\cdot}{p}\bigr)$.
Define $f(pa)=\chi_p(a)$ for $p\nmid a$, $f(n)=-\chi_p(n)/p$ for $p\nmid n$, and $f(0)=0$.
Then $f$ is an odd Okada vector for $p^2$; in fact $L(1,f)=0$.
\end{proposition}

\begin{proof}
The values $\pm1$ and $\mp1/p$ are odd elements of $\Zt$, and $\sum_{n=1}^{p^2}f(n)=0$ because $\chi_p$ sums to $0$ over any complete set of units.
Moreover
\[
L(1,f)=\sum_{p\nmid a}\frac{\chi_p(a)}{pa}-\frac1p\sum_{p\nmid n}\frac{\chi_p(n)}{n}=\frac{L(1,\chi_p)}p-\frac{L(1,\chi_p)}p=0. \qedhere
\]
\end{proof}

\begin{theorem}\label{thm:squares}
Every odd square $q>1$ has an odd Okada vector.
\end{theorem}

\begin{proof}
Write $q=p_1^2p_2^2\cdots p_r^2$ with primes $p_i$ (repetitions allowed).
Lemma~\ref{lem:compose} applies verbatim to odd Okada vectors: if $f\in\Ok_q$ and $g\in\Ok_N$ are odd and $2$-integral, the function $F$ of the lemma has $F(n)\in\{f(n),g(n/q)\}$ odd and $2$-integral for $qN\nmid n$, $\sum F=0$, and $L(1,F)=L(1,f)+L(1,g)/q=0$, so $F\in\Ok_{qN}$ by Theorem~\ref{thm:okada}.
Starting from Proposition~\ref{prop:p2} and composing $r-1$ times gives an odd Okada vector for $q$.
\end{proof}

Thus the $2$-adic method can never exclude a square modulus, and squares are the only place where it can fail to prove Conjecture~\ref{conj:erdos} if the following conjecture holds.

\subsection{Non-squares}
\begin{conjecture}\label{conj:nonsquare}
Every odd non-square $q$ is $2$-adically obstructed.
Equivalently: an odd $q>1$ has an odd Okada vector if and only if $q$ is a perfect square.
\end{conjecture}

Conjecture~\ref{conj:nonsquare} implies Conjecture~\ref{conj:erdos} for all non-square $q$, so it would confine all counterexamples to square moduli.
It holds for odd prime powers $q=p^{e}$ with $e$ odd: the layers are $d=p^j$ ($1\le j\le e$), all with $Q_d=C_{2^{t_p}}$, the only conductor is $m=p$, and (A$'$) says that $\Gamma=\sum_{j}\Phi_{p^j}$ has vanishing non-trivial Fourier coefficients, i.e.\ is constant; if all $\Phi_{p^j}(y)$ are odd, the constant is a sum of $e$ odd numbers, hence odd, and then $\sum_dS_d=2^{t_p}\cdot(\text{odd})\neq0$ contradicts (C$'$).
For $q\equiv3\pmod4$ it is Proposition~\ref{prop:MS}.
The shadow system also yields a residuacity criterion that covers infinitely many new moduli $q\equiv1\pmod 4$.

\begin{theorem}[Residuacity criterion]\label{thm:resid}
Let $q$ be odd and let $p$ be a prime with $p\parallel q$.
Put $t=v_2(p-1)$.
If every prime $r\neq p$ dividing $q$ is a $2^{t}$-th power residue modulo $p$, then $q$ is $2$-adically obstructed; in particular Conjecture~\ref{conj:erdos} holds for $q$.
\end{theorem}

\begin{proof}
Suppose $\Phi$ is an odd $2$-integral solution of the shadow system, and write $\Phi_d=|H_d|\mathbf 1-2B_d$ with $B_d\colon Q_d\to\Zt$ (possible since $\Phi_d$ and $|H_d|$ are odd).
Let $\chi$ be a non-trivial character of $Q_p=C_{2^t}$, viewed as a primitive character of conductor $p$.
By hypothesis $\chi(r)=1$ for every prime $r\mid q$, $r\ne p$, so in (A$'$) every weight equals $\frac{d}{\varphi(d)}\prod_{r\mid d,\,r\neq p}(1-\frac1r)=\frac p{p-1}$, and (A$'$) becomes $\sum_{p\mid d}\sum_y\Phi_d(y)\overline{\chi(y)}=0$, i.e.\ $\sum_{p\mid d}\ip{B_d}{\chi}=0$.
Hence the function $\Gamma=\sum_{p\mid d}\pi_{d\to p}B_d$ on $C_{2^t}$ has vanishing non-trivial Fourier coefficients, so it is a constant $c$, and $c=\Gamma(0)\in\Zt$.
Therefore
\[
\sum_{p\mid d}\sum_{y}B_d(y)=2^{t}c\equiv0\pmod{2^t}.
\]
On the other hand, since $p\parallel q$ we have $\varepsilon_p(d)=1$ for $p\nmid d$ and $\varepsilon_p(d)=-1/(p-1)$ for $p\mid d$, so (B$'$) and (C$'$) give $\sum_{p\mid d}S_d=0$, where $S_d=\sum_y\Phi_d(y)=\varphi(d)-2\sum_yB_d(y)$.
Since $\sum_{p\mid d}\varphi(d)=q-q/p$, this yields $\sum_{p\mid d}\sum_yB_d(y)=\frac qp\cdot\frac{p-1}2$, whose $2$-adic valuation is $t-1$, a contradiction.
\end{proof}

\begin{corollary}\label{cor:resid}
There are infinitely many $q\equiv1\pmod4$ with $2\varphi(q)+1\le q$ (so that Okada's criterion does not apply) for which Conjecture~\ref{conj:erdos} holds; for instance $q=3\cdot7\cdot13\cdot19\cdot31=160797$.
\end{corollary}

\begin{proof}
Take $p=3$ and $q=3\,r_1\cdots r_k$ with distinct primes $r_i\equiv1\pmod 3$; every $r_i$ is then a square modulo $3$.
Because $\sum_{r\equiv1\,(3)}1/r$ diverges, $\varphi(q)/q<\frac12$ once $k$ is large, and the condition $q\equiv1\pmod4$ can be met by including or omitting one prime $r\equiv7\pmod{12}$.
For $160797$ one has $\varphi(q)/q\approx0.484$; this modulus is also not covered by \cite[Prop.~3.1, Cor.~3.3]{CM15} or \cite[Cor.~1]{Bha22}.
\end{proof}

Among the $74$ new moduli of Remark~\ref{rem:new} below, Theorem~\ref{thm:resid} applies (with $p=11$) to $165$, $825$, $1485$, $4125$, $6105$ and $7425$; for the others, such as $105$ and $945$, no single prime suffices and we rely on computation.
We do not have a proof of Conjecture~\ref{conj:nonsquare} in general; the evidence is the following computation.

\begin{theorem}[computer-assisted]\label{thm:nonsquares}
Every odd non-square $q\le 13105$ is $2$-adically obstructed.
\end{theorem}

\begin{proof}
For primes this is Baker--Birch--Wirsing (or the case $e=1$ above).
For each of the $4938$ odd composite non-squares $q\le13105$ we computed a certificate in the sense of Lemma~\ref{lem:farkas} for the shadow system of Theorem~\ref{thm:shadow} in the integer form of Remark~\ref{rem:integer-form}: a vector $\Lambda$ with $2$-integral rational entries, one for each unknown $\Phi_d(y)$, which lies in the $\Q$-span of the equations and satisfies $\Lambda(\mathbf1)\not\equiv0\pmod 2$.
If $\Phi$ were an odd $2$-integral solution, then $0=\Lambda(\Phi)\equiv\Lambda(\mathbf1)\equiv1\pmod2$.
The certificates and an independent checker are part of the supplementary material; the checker verifies $2$-integrality and oddness directly and membership in the row space by an exact rank computation over $\Q$ (Appendix~\ref{app:cert}).
\end{proof}

The same computation confirms Theorem~\ref{thm:squares} in this range: for each of the $56$ odd squares $9\le q\le13105$ the shadow system has odd solutions.
As cross-checks, the descent was compared with the full system of Theorem~\ref{thm:okada} in all $q-1$ unknowns for $39$ moduli up to $2145$ (Section~\ref{sec:found}); and a second program, written independently from the statement of Okada's original equations (it uses neither characters, nor layers, nor the descent), decided the existence of odd Okada vectors for all $332$ odd composite $q\le999$ by two methods ($2$-adic elimination over $\Zt$, and exact kernel saturation): they exist exactly for the $15$ odd squares $3^2,5^2,\dots,31^2$, in agreement with Theorems~\ref{thm:squares} and~\ref{thm:nonsquares}.

\begin{corollary}\label{cor:transc}
For every odd non-square $q\le13105$ and every Erd\H{o}s function $f$ modulo $q$ with $\sum_{n=1}^qf(n)=0$, the number $\sum_{n\ge1}f(n)/n$ is non-zero, hence transcendental.
\end{corollary}

\begin{proof}
Theorem~\ref{thm:nonsquares} and Lemma~\ref{lem:erdos-odd}.
\end{proof}

\begin{remark}\label{rem:new}
Of the moduli in Corollary~\ref{cor:transc}, exactly $74$ satisfy none of the sufficient conditions known to us from the literature: $q$ prime \cite{BBW73}; $2\varphi(q)+1>q$ \cite{Oka82}; $q\equiv3\pmod4$ \cite{MS10}; $\sum_{d\mid q,\,d\ge3}1/\varphi(d)<1$ or ``least prime factor $\ge$ number of divisors'' \cite[Prop.~3.1, Cor.~3.3]{CM15}; and the family $q=pN$ of \cite[Cor.~1]{Bha22}.
They are all $\equiv1\pmod4$; the first ones are
$105$, $165$, $525$, $585$, $825$, $945$, $1365$, $1485$, $1785$, $2145$, $2205$, $2625$, $2805$, $2925$, $3045$, $3465$.
\end{remark}

\subsection{The least counterexample}
\begin{theorem}\label{thm:least}
Let $2\le q\le13105$.
An Erd\H{o}s function modulo $q$ with $\sum_{n=1}^qf(n)=0$ and $L(1,f)=0$ exists if and only if $q=11025$.
\end{theorem}

\begin{proof}
If $q$ is even the sum $\sum_{n=1}^qf(n)$ has an odd number of terms $\pm1$ and cannot vanish.
If $q$ is an odd non-square, use Theorem~\ref{thm:nonsquares} and Lemma~\ref{lem:erdos-odd}.
If $q=m^2$ is an odd square, then $m\le114$, and every odd $m\le114$ other than $105=3\cdot5\cdot7$ has at most two distinct prime factors (the next candidate is $3\cdot5\cdot11=165$); hence $\varphi(q)/q=\prod_{p\mid m}(1-1/p)\ge\frac23\cdot\frac45>\frac12$ and Okada's theorem \cite{Oka82} applies.
Finally $q=105^2$ is Theorem~\ref{thm:11025}.
\end{proof}

\section{How the counterexamples were found, and what failed}\label{sec:found}

\subsection{Symmetric lifting}
By Theorem~\ref{thm:shadow}(2), an $H$-invariant Erd\H{o}s function modulo $q$ is the same thing as a solution of the shadow system with $\Phi_d(y)=|H_d|\,s_d(y)$, $s_d(y)\in\{\pm1\}$.
For a square $q$ the shadow system does have odd solutions (Theorem~\ref{thm:squares}), so it is reasonable to look for $\pm1$ solutions among them.
This is a $0/1$ feasibility problem with $\prod(1+e_p2^{t_p})-1$ variables and the integer equations of Remark~\ref{rem:integer-form}; for $q=105^2$ and $q=165^2$ there are $224$ variables and $19$ equations.
We solved it with the CP-SAT solver of OR-Tools, version 9.15 \cite{ORT}.
The outcomes, with the time limits used, were as follows.
\begin{itemize}
\item $q=105^2$: a solution was found in about $214$ seconds (two earlier runs with limits of $200$ and $600$ seconds ended without a decision).
Minimizing the number of shadow classes on which $s_d$ differs from $\pm\chi_{35}$ on the layers $35\mid d$ then produced, within $900$ seconds, the solution of Table~\ref{tab:11025}, which differs from that pattern only on the layer $d=105$ and depends on Legendre symbols alone (optimality of this count was not proved).
\item $q=165^2$: a solution was found in about $55$ seconds.
The subproblems ``$f$ depends only on Legendre symbols'', ``$f$ odd'' ($f(-n)=-f(n)$) and ``$f$ even'' ($f(-n)=f(n)$) are all infeasible (proved by the solver).
\item $q=195^2=38025$: the $H$-invariant problem is infeasible (proved in about one second), so a counterexample modulo $195^2$, if any, cannot be $H$-invariant.
\item Legendre-type $H$-invariant problems (each $s_d$ a function of the Legendre symbols only): solutions were found for $315^2$ (in under one second), $735^2$ (about $67$ seconds) and $495^2$ (about $75$ seconds), giving Theorem~\ref{thm:more}; the problem is infeasible for $165^2$ and $195^2$, and was undecided for $525^2$ after $600$ seconds.
\item The squares $225, 441, 2025, 53361, 65025,\dots$ with $2\varphi(q)+1>q$ gave infeasible problems, as they must by Okada's theorem.
\item Functions depending only on $\gcd(n,q)$ never give counterexamples for the squares we tried ($15^2,105^2,165^2,195^2,315^2,1155^2$): the corresponding sign problems, which only involve the conditions (B$'$) and (C$'$), are infeasible.
\end{itemize}
Each solution was expanded to all $q-1$ residues and verified exactly as described in Section~\ref{sec:counter}.

\subsection{Relaxations that do not suffice}
The descent of Theorem~\ref{thm:shadow} keeps the whole maximal $2$-power quotient.
Several natural smaller systems are too weak to detect the obstruction for non-squares:
keeping only the quadratic characters (in the form valid for relaxations: unknowns are the fibre sums of $b=(1-f)/2$, and (B$'$), (C$'$) become inhomogeneous) misses the obstruction for $116$ of the $248$ odd composite non-squares $q<800$, the first being $15,35,45$; capping all $t_p$ at $2$ still misses $40$ of them, the first being $51=3\cdot17$;
keeping only characters supported on one prime misses $15,35,45,51,\dots$, and on at most two primes misses $105,195,231,255,945,1155,1365$ among the moduli tested;
keeping only the odd characters detects exactly the $q\equiv3\pmod4$ among the $19$ moduli tested (as Proposition~\ref{prop:MS} predicts), and keeping only the even characters detects nothing.
So for $q\equiv1\pmod4$ the obstruction genuinely involves the interaction of odd and even characters, of all $2$-power orders, supported on all primes of $q$ at once.

\subsection{Validation of the computations}
The layered system of Theorem~\ref{thm:okada} was validated before use: for $q\le315$ (fifteen moduli) its rank is $\varphi(q)+\omega(q)$, random rational vectors in its kernel give $|L(1,f)|<10^{-50}$ in $50$-digit arithmetic, and generic vectors do not.
The shadow test of Theorem~\ref{thm:shadow} agrees with the test on the full system in all $q-1$ unknowns for the $39$ moduli we compared (all $26$ odd composite $q\le105$ and thirteen larger ones up to $2145$), and the two forms of the shadow test (unknowns $\Phi$ with target ``all odd'', and unknowns $B$ with inhomogeneous right-hand side) agree for all odd composite $q<800$.
The independent program mentioned after Theorem~\ref{thm:nonsquares} also re-verified the counterexample of Table~\ref{tab:11025}, the first counterexample found modulo $105^2$ (it is in the supplement), and the counterexamples modulo $165^2$ and $315^2$: all of Okada's equations hold exactly, in two different implementations, and $L(1,f)$ vanishes to more than $110$ digits by three numerical methods, one of which (exact partial sums up to $40q$ plus a moment expansion of the tail in Hurwitz zeta values) does not use the digamma function; perturbing $f$ at two places makes $L(1,f)$ visibly non-zero.

\section{Open problems}\label{sec:open}

\begin{enumerate}
\item \emph{Prove Conjecture~\ref{conj:nonsquare}.}
It is known for $q\equiv3\pmod 4$ and for odd prime powers with odd exponent, and verified for $q\le13105$.
A proof would establish Conjecture~\ref{conj:erdos} for all non-square $q$, together with the transcendence of all the corresponding sums.
The relaxation experiments of Section~\ref{sec:found} indicate that a proof must use the full $2$-power quotient and the interaction of odd and even characters.
A natural intermediate statement would be a lifting lemma: if $q$ is obstructed and $\gcd(q,N)=1$, then $qN$ is obstructed; together with the prime-power case it would imply the conjecture.
\item \emph{Determine the set $\mathcal C$ of counterexample moduli.}
We know $105^2,165^2,315^2,495^2,735^2\in\mathcal C$, that $\mathcal C$ is a multiplicative semigroup, that every element of $\mathcal C$ is odd, and that $q\notin\mathcal C$ whenever $2\varphi(q)+1>q$.
If Conjecture~\ref{conj:nonsquare} holds, $\mathcal C$ is contained in the set of odd squares $q$ with $2\varphi(q)+1\le q$.
Is $195^2\in\mathcal C$? (No counterexample modulo $195^2$ can be invariant under the odd-order units.)
Does every odd square $q$ with $2\varphi(q)+1\le q$ belong to $\mathcal C$?
\item \emph{Explain the counterexamples.}
The function of Table~\ref{tab:11025} is built from the Jacobi symbol $\bigl(\frac{x}{35}\bigr)$ on most layers.
Is there a conceptual construction, for instance through Gauss sums or class numbers of the quadratic fields of conductor dividing $105$, that produces counterexamples for an infinite family of squares beyond the semigroup generated by $105^2$ and $165^2$?
\item \emph{Higher $L$-values.}
The analogous questions for $\sum f(n)/n^k$, $k\ge2$, are tied to the Chowla--Milnor conjecture \cite{Bha22}; it would be interesting to know whether a $2$-adic shadow principle also applies there.
\end{enumerate}

\appendix
\section{\texorpdfstring{The counterexample modulo $165^2$}{The counterexample modulo 165^2}}\label{app:27225}
\begin{longtable}{rrrl}
\caption{The counterexample modulo $q=27225=3^2 5^2 11^2$. For $n=(q/d)x$ with $\gcd(x,d)=1$ put $\alpha=\operatorname{ind}_2(x\bmod 3)\bmod 2$, $\beta=\operatorname{ind}_2(x\bmod 5)\bmod 4$, $\gamma=\operatorname{ind}_2(x\bmod 11)\bmod 2$ (only for the primes dividing $d$; $2$ is a primitive root modulo $3,5,11$). The string lists $f(n)$ for the tuples $(\alpha,\beta,\gamma)$ in lexicographic order (coordinates for primes not dividing $d$ omitted).}\label{tab:27225}\\
\toprule $d$ & $q/d$ & coordinates & signs\\ \midrule\endfirsthead
\toprule $d$ & $q/d$ & coordinates & signs\\ \midrule\endhead
3 & 9075 & $(\alpha)$ & \texttt{++}\\
5 & 5445 & $(\beta)$ & \texttt{----}\\
9 & 3025 & $(\alpha)$ & \texttt{++}\\
11 & 2475 & $(\gamma)$ & \texttt{-+}\\
15 & 1815 & $(\alpha,\beta)$ & \texttt{--+++---}\\
25 & 1089 & $(\beta)$ & \texttt{+-+-}\\
33 & 825 & $(\alpha,\gamma)$ & \texttt{+-+-}\\
45 & 605 & $(\alpha,\beta)$ & \texttt{-+---+--}\\
55 & 495 & $(\beta,\gamma)$ & \texttt{+--++--+}\\
75 & 363 & $(\alpha,\beta)$ & \texttt{++--+-++}\\
99 & 275 & $(\alpha,\gamma)$ & \texttt{--++}\\
121 & 225 & $(\gamma)$ & \texttt{-+}\\
165 & 165 & $(\alpha,\beta,\gamma)$ & \texttt{+--+----+--+++++}\\
225 & 121 & $(\alpha,\beta)$ & \texttt{-+-++--+}\\
275 & 99 & $(\beta,\gamma)$ & \texttt{-+-+++--}\\
363 & 75 & $(\alpha,\gamma)$ & \texttt{++--}\\
495 & 55 & $(\alpha,\beta,\gamma)$ & \texttt{-+---+-+-+++-+-+}\\
605 & 45 & $(\beta,\gamma)$ & \texttt{-+-+-+-+}\\
825 & 33 & $(\alpha,\beta,\gamma)$ & \texttt{-+-+-+-+-+-+-+-+}\\
1089 & 25 & $(\alpha,\gamma)$ & \texttt{-+-+}\\
1815 & 15 & $(\alpha,\beta,\gamma)$ & \texttt{-+-+-+-+-+-+-+-+}\\
2475 & 11 & $(\alpha,\beta,\gamma)$ & \texttt{-+-+-+-+-+-+-+-+}\\
3025 & 9 & $(\beta,\gamma)$ & \texttt{-+-+-+-+}\\
5445 & 5 & $(\alpha,\beta,\gamma)$ & \texttt{-+-+-+-+-+-+-+-+}\\
9075 & 3 & $(\alpha,\beta,\gamma)$ & \texttt{-+-+-+-+-+-+-+-+}\\
27225 & 1 & $(\alpha,\beta,\gamma)$ & \texttt{+-+-+-+-+-+-+-+-}\\
\bottomrule\end{longtable}


\section{Certificates and algorithms}\label{app:cert}

\subsection*{Deciding obstruction}
For a modulus $q$ we build the integer system of Remark~\ref{rem:integer-form}: rows indexed by the primitive functionals of the $H_m(\Phi)$ for squarefree $m\mid\rad(q)$, $m>1$, the rows (B$'$) multiplied by $p-1$, and the row (C$'$).
We compute a basis of its rational kernel, saturate it at $2$ (repeatedly replacing a basis vector by half of an $\F_2$-dependent combination until the reductions modulo $2$ are independent), and test whether the all-ones vector lies in the $\F_2$-span of the reductions.
By Lemma~\ref{lem:farkas} this decides $2$-adic obstruction.

\subsection*{Certificates}
If $q$ is obstructed, let $k_1,\dots,k_r$ be the saturated kernel basis.
Solving the $\F_2$-system $\bar\Lambda\cdot k_j=0$ $(1\le j\le r)$, $\bar\Lambda\cdot\mathbf1=1$ gives $\bar\Lambda\in\{0,1\}^N$; then $\Lambda=\bar\Lambda+2w$, where $w$ is supported on the pivot coordinates of the $k_j$ and solves $w\cdot k_j=-\frac12\bar\Lambda\cdot k_j$, satisfies $\Lambda\cdot k_j=0$ for all $j$ (so $\Lambda$ lies in the row space), has $2$-integral entries (the pivot minor is odd), and $\Lambda(\mathbf1)\equiv\bar\Lambda(\mathbf1)=1\pmod 2$.
The supplementary file of certificates contains $\Lambda$ for each of the $4938$ odd composite non-squares $q\le 13105$.
The checker rebuilds the integer system from its definition and verifies, for each $q$: every entry of $\Lambda$ has odd denominator; $\Lambda(\mathbf1)$ has odd numerator and denominator; and $\operatorname{rank}\binom{M}{\Lambda}=\operatorname{rank}M$ over $\Q$, computed exactly with FLINT \cite{FLINT}.
These checks do not depend on the saturation step.

\subsection*{Verification of the counterexamples}
For the two counterexamples we verified, independently of the shadow code: the values are $\pm1$ off multiples of $q$ and sum to $0$; Okada's original equations hold exactly (all $\varphi(q)$ unit equations, with the infinite sums over $M(q)$ evaluated as geometric series, and the $\omega(q)$ prime equations); the coefficients \eqref{eq:A} vanish for all characters of $2$-power order (for $105^2$ by hand, Tables~\ref{tab:chars}--\ref{tab:logs}); and $L(1,f)=-\frac1q\sum_af(a)\psi(a/q)$ vanishes to $134$ significant digits in PARI/GP and to $50$ digits in mpmath.

\subsection*{The $74$ new moduli of Remark~\ref{rem:new}}
$105$, $165$, $525$, $585$, $825$, $945$, $1365$, $1485$, $1785$, $2145$, $2205$, $2625$, $2805$, $2925$, $3045$, $3465$, $3705$, $3885$, $4125$, $4305$, $4389$, $4485$, $4641$, $4725$, $4785$, $4845$, $5145$, $5265$, $5313$, $5445$, $5565$, $5865$, $5985$, $6045$, $6105$, $6405$, $6765$, $6825$, $7245$, $7425$, $7605$, $7665$, $7905$, $8085$, $8265$, $8385$, $8505$, $8745$, $8925$, $9009$, $9165$, $9345$, $9405$, $9765$, $9945$, $10005$, $10065$, $10185$, $10545$, $10605$, $10725$, $10965$, $11385$, $11445$, $11505$, $11685$, $11781$, $11865$, $11985$, $12045$, $12285$, $12705$, $12765$, $13065$.

\section{Code and data}\label{app:code}
The supplementary archive contains: the layered criterion and its validation; the shadow construction and the obstruction test in both forms; the certificate generator and checker with the certificate file for $q\le13105$; the CP-SAT models used to find the counterexamples; the two counterexamples as explicit lists of values $f(1),\dots,f(q-1)$; the independent verification scripts (exact Okada equations, character coefficients, PARI/GP numerics); and logs of all runs quoted in the text.

\begin{thebibliography}{99}
\bibitem{Bak75} A. Baker, \emph{Transcendental Number Theory}, Cambridge University Press, 1975.
\bibitem{BBW73} A. Baker, B. J. Birch, E. A. Wirsing, On a problem of Chowla, \emph{J. Number Theory} \textbf{5} (1973), 224--236.
\bibitem{Bha22} A. Bharadwaj, On primitivity and vanishing of Dirichlet series, preprint, arXiv:2204.03674 (2022).
\bibitem{CM15} T. Chatterjee, M. R. Murty, On a conjecture of Erd\H{o}s and certain Dirichlet series, \emph{Pacific J. Math.} \textbf{275} (2015), no.~1, 103--113; arXiv:1501.04185.
\bibitem{Cho70} S. Chowla, The nonexistence of nontrivial linear relations between the roots of a certain irreducible equation, \emph{J. Number Theory} \textbf{2} (1970), 120--123.
\bibitem{Liv65} A. E. Livingston, The series $\sum_1^\infty f(n)/n$ for periodic $f$, \emph{Canad. Math. Bull.} \textbf{8} (1965), 413--432.
\bibitem{MP22} M. R. Murty, S. Pathak, The Okada space and vanishing of $L(1,f)$, \emph{Funct. Approx. Comment. Math.} \textbf{66} (2022), no.~1; doi:10.7169/facm/1952.
\bibitem{MS07} M. R. Murty, N. Saradha, Transcendental values of the digamma function, \emph{J. Number Theory} \textbf{125} (2007), 298--318.
\bibitem{MS10} M. R. Murty, N. Saradha, Euler--Lehmer constants and a conjecture of Erd\H{o}s, \emph{J. Number Theory} \textbf{130} (2010), 2671--2682.
\bibitem{Oka82} T. Okada, On a certain infinite series for a periodic arithmetical function, \emph{Acta Arith.} \textbf{40} (1982), 143--153.
\bibitem{Pat17} S. Pathak, On a conjecture of Livingston, \emph{Canad. Math. Bull.} \textbf{60} (2017), 184--195.
\bibitem{Pat19} S. Pathak, Distribution and non-vanishing of special values of $L$-series attached to Erd\H{o}s functions, \emph{Int. J. Number Theory} \textbf{15} (2019); doi:10.1142/S1793042119500829; arXiv:1909.03972.
\bibitem{ST03} N. Saradha, R. Tijdeman, On the transcendence of infinite sums of values of rational functions, \emph{J. London Math. Soc. (2)} \textbf{67} (2003), 580--592.
\bibitem{ORT} L. Perron, F. Didier, \emph{CP-SAT}, OR-Tools 9.15, Google, \url{https://developers.google.com/optimization}.
\bibitem{FLINT} The FLINT team, \emph{FLINT: Fast Library for Number Theory}, \url{https://flintlib.org}; python-flint 0.9.
\bibitem{PARI} The PARI~Group, \emph{PARI/GP}, Univ. Bordeaux, \url{https://pari.math.u-bordeaux.fr}.
\end{thebibliography}
\end{document}

